\documentclass[10pt,a4paper]{article}
\usepackage{amsmath,amssymb,amsthm,graphicx}
\usepackage{hyperref}

\usepackage[scr=rsfs,cal=euler]{mathalfa}
\usepackage{bbm}
\let\mathds\mathbbm
\multlinegap0pt

\newenvironment{enumeratei}
{\bgroup\def\theenumi{\roman{enumi}}\def\theenumii{\arabic{enumii}}\begin{enumerate}}
{\end{enumerate}\egroup}
\newenvironment{enumerateI}
{\bgroup\def\theenumi{\Roman{enumi}}\def\theenumii{\arabic{enumii}}\begin{enumerate}}
{\end{enumerate}\egroup}

\newcommand{\Psfrac}[2]{(\sfrac{#1}{#2})}
\newcommand{\psfrac}[2]{\sfrac{(#1)}{#2}}
\newcommand{\spfrac}[2]{\sfrac{#1}{(#2)}}
\newcommand{\pspfrac}[2]{\sfrac{(#1)}{(#2)}}
\newcommand\mto{\mathchoice{\longmapsto}{\mapsto}{\mapsto}{\mapsto}}

\usepackage{mathtools}
\usepackage{esint}

\theoremstyle{plain}
\newtheorem{theor10}{Theorem}
\newenvironment{theor1}
{\renewcommand{\qedsymbol}{$\lozenge$}\pushQED{\qed}\begin{theor10}}
{\popQED\end{theor10}}
\newtheorem{cor10}{Corollary}
\newenvironment{cor1}
{\renewcommand{\qedsymbol}{$\lozenge$}\pushQED{\qed}\begin{cor10}}
{\popQED\end{cor10}}
\newtheorem{theor0}{Theorem}[section]
\newenvironment{theor}
{\renewcommand{\qedsymbol}{$\lozenge$}\pushQED{\qed}\begin{theor0}}
{\popQED\end{theor0}}
\newtheorem{lem0}[theor0]{Lemma}
\newenvironment{lem}
{\renewcommand{\qedsymbol}{$\lozenge$}\pushQED{\qed}\begin{lem0}}
{\popQED\end{lem0}}
\newtheorem{prop0}[theor0]{Proposition}
\newenvironment{prop}
{\renewcommand{\qedsymbol}{$\lozenge$}\pushQED{\qed}\begin{prop0}}
{\popQED\end{prop0}}
\newtheorem{cor0}[theor0]{Corollary}
\newenvironment{cor}
{\renewcommand{\qedsymbol}{$\lozenge$}\pushQED{\qed}\begin{cor0}}
{\popQED\end{cor0}}
\theoremstyle{definition}
\newtheorem{defin0}[theor0]{Definition}
\newenvironment{defin}
{\renewcommand{\qedsymbol}{$\lozenge$}\pushQED{\qed}\begin{defin0}}
{\popQED\end{defin0}}
\newtheorem{rems0}[theor0]{Remarks}
\newenvironment{rems}
{\renewcommand{\qedsymbol}{$\lozenge$}\pushQED{\qed}\begin{rems0}}
{\popQED\end{rems0}}
\newtheorem{rem0}[theor0]{Remark}
\newenvironment{rem}
{\renewcommand{\qedsymbol}{$\lozenge$}\pushQED{\qed}\begin{rem0}}
{\popQED\end{rem0}}

\numberwithin{equation}{section}

\newcommand{\N}{\mathbb N}
\newcommand{\e}{\varepsilon}
\newcommand{\Lc}{\mathcal{L}}
\newcommand{\Vc}{\mathcal V}
\newcommand{\R}{\mathbb R}
\newcommand{\E}{\mathbb E}
\newcommand{\C}{\mathbb C}
\newcommand{\F}{\mathcal F}
\newcommand{\loc}{{\operatorname{loc}}}
\newcommand{\Id}{\operatorname{Id}}
\newcommand{\per}{{\operatorname{per}}}
\newcommand{\ee}{e}
\newcommand{\Aa}{\boldsymbol a}
\newcommand{\bb}{{\boldsymbol b}}
\newcommand{\Bb}{{\boldsymbol B}}
\newcommand{\aA}{{\boldsymbol A}}
\newcommand{\cc}{{\boldsymbol c}}
\newcommand{\Ld}{\operatorname{L}}
\newcommand{\supp}{\operatorname{supp}}
\newcommand{\diam}{\operatorname{diam}}
\newcommand{\step}[1]{\subsubsection*{Step \textup #1}}
\newcommand{\Step}[1]{\Subsubsection*{Step \textup #1}}
\newcommand{\Pm}{\mathbb{P}}
\newcommand{\expec}[1]{\mathbb{E}\left[ #1 \right]}
\newcommand{\Bexpec}[1]{\mathbb{E}\Bigl[ #1 \Bigr]}
\newcommand{\expecm}[1]{\mathbb{E}\bigl[ #1 \bigr]}\let\bexpec\expecm

\title{A priori estimates for the wave equation with a self-adjoint operator squeezed between two Laplacians}
\author{Mitia Duerinckx, Antoine Gloria, Matthias Ruf}
\date{doi:10.5802/jep.259; CC BY 4.0}
\begin{document}
\maketitle
\setcounter{section}{1}
\noindent\textit{Editorial note.} All delivered mathematical text is the author's own, copied from the frozen source sources/jep-2024-259-a-spectral-ansatz-for-the-long-time-homogenization/pinned-original-source.tex; no statement or proof step was rewritten, reordered or removed.\par\medskip
\noindent\textit{Editorial note.} The unit has no row in the census, so the gate is the only authority for this card, printed numbers are computed by hand and the finished PDF is verified with pdftotext (duty desk rule for this band).\par\medskip
\noindent\textit{Editorial note.} The counted span carries no \textbackslash{}cite and no reference to a label outside the span.\par\medskip
\noindent\textit{Editorial scope.} Single Lemma card cut from doi:10.5802/jep.259 (A spectral ansatz for the long-time homogenization of the wave equation); statement at source lines 2383-2397 and the author's own proof at source lines 2398-2414.\par\medskip
\noindent\textit{Allowed standard prerequisites (editorial).}\begin{itemize}
\item The two-sided bound on the operator. The hypothesis of the statement, source line 2384. It makes the square root of the operator a bounded operator from the space into the homogeneous order minus one space with bounded inverse, which is what allows the divergence in the source term to be absorbed.
\item Duhamel's formula for the wave equation. The representation used in the proof at source lines 2401-2406. It writes the solution as an integral of the sine propagator against the source terms, which is then estimated by spectral calculus.
\end{itemize}\medskip
\begin{lem}[A priori estimates]\label{lem:apriori}
Let $\Lc$ be a self-adjoint operator on $\Ld^2(\R^d)$ satisfying the bound $-\triangle\le\Lc\le-C_0\triangle$ for some constant $C_0<\infty$. Given
$F_1\in\Ld^1_\loc(\R^+;\Ld^2(\R^d))$ and $F_2,F_3\in W^{1,1}_\loc(\R^+;\Ld^2(\R^d))$ with $F_2|_{t=0}=F_3|_{t=0}=0$, let $z$ be the solution of the wave equation
\begin{equation*}
\begin{cases}
(\partial_t^2+\Lc) z=F_1+\nabla \cdot F_2+\partial_t F_3\quad\text{in $\R^d$},\\
z|_{t=0}=\partial_t z|_{t=0}=0.
\end{cases}
\end{equation*}
Then, for all $t\ge 0$, we have
\[\|z^t\|_{\Ld^2(\R^d)}\,\lesssim_{C_0}\,t\|F_1\|_{\Ld^1((0,t);\Ld^2(\R^d))}+\|(F_2,F_3)\|_{\Ld^1((0,t);\Ld^2(\R^d))},\]
and
\[\|Dz^t\|_{\Ld^2(\R^d)} \,\lesssim_{C_0}\,\|(F_1,\partial_tF_2,\partial_tF_3)\|_{\Ld^1((0,t);\Ld^2(\R^d))}.\qedhere\]
\end{lem}


\begin{proof}
The assumption $-\triangle\le\Lc\le-C_0\triangle$ entails that $\sqrt\Lc$ defines a bounded linear operator $\Ld^2(\R^d)\to\dot H^{-1}(\R^d)$ with bounded inverse. We may therefore define~$\tilde F_2$ as the solution of $\sqrt{\Lc} \tilde F_2=\nabla\cdot F_2$, which satisfies $\|\tilde F_2\|_{\Ld^2(\R^d)}\lesssim_{C_0}\|F_2\|_{\Ld^2(\R^d)}$.
The solution $z$ of the wave equation can then be represented in terms of Duhamel's formula,
\[z^t=\int_0^t\frac{\sin((t-s)\sqrt\Lc)}{\sqrt\Lc}\,(F_1^s+\partial_s F_3^s)\,ds+\int_0^t\sin((t-s)\sqrt\Lc)\,\tilde F_2^s\,ds.\]
Integrating by parts in the integral for $F_3$, with $F_3|_{t=0}=0$, this can be rewritten as
\[z^t=\int_0^t\frac{\sin((t-s)\sqrt\Lc)}{\sqrt\Lc}\,F_1^s\,ds+\int_0^t\sin((t-s)\sqrt\Lc)\,\tilde F_2^s\,ds+\int_0^t\cos((t-s)\sqrt\Lc)\,F_3^s\,ds,\]
and the stated $\Ld^2$-estimate follows from spectral calculus.
For the energy estimate, we rather use energy conservation in form of
\begin{align*}
\tfrac12\partial_t\int_{\R^d}\big(|\partial_tz|^2+z\Lc z+2F_2\cdot\nabla z\big)
&=\int_{\R^d} (\partial_tz)\,(\partial^2_t+\Lc)z+\partial_t\int_{\R^d}F_2\cdot\nabla z\\
&=\int_{\R^d}(\partial_tz)\big(F_1+\partial_tF_3+\nabla\cdot F_2\big)+\partial_t\int_{\R^d}F_2\cdot\nabla z\\
&=\int_{\R^d}(\partial_tz)\big(F_1+\partial_tF_3\big)+\int_{\R^d}\partial_tF_2\cdot\nabla z\\
&\le\|Dz\|_{\Ld^2(\R^d)}\|(F_1,\partial_tF_2,\partial_t F_3)\|_{\Ld^2(\R^d)},
\end{align*}
and the conclusion follows.
\end{proof}

\end{document}
